\documentclass[11pt]{article}
\usepackage[utf8]{inputenc}
\usepackage[T1]{fontenc}
\usepackage{lmodern}
\usepackage[english]{babel}
\usepackage[margin=1in]{geometry}
\usepackage{titlesec}
\usepackage{enumitem}
\usepackage{hyperref}
\usepackage{parskip}

\hypersetup{
  colorlinks=true,
  urlcolor=black,
  linkcolor=black,
}

\pagestyle{plain}
\setlength{\parindent}{0pt}
\setlength{\parskip}{0.35em}

\titleformat{\section}
  {\normalsize\bfseries\MakeUppercase}
  {}{0em}{}[\vspace{-0.35em}\hrule\vspace{0.85em}]

\titlespacing*{\section}{0pt}{1.6em}{1.0em}

\newcommand{\cvsubsection}[1]{\vspace{0.9em}\textit{#1}\par\vspace{0.35em}}

\newcommand{\entry}[2]{\noindent\textbf{#1}\hfill #2\par\vspace{0.15em}}
\newcommand{\entryplain}[2]{\noindent#1\hfill #2\par\vspace{0.15em}}
\newcommand{\detail}[1]{\noindent\hspace*{1.5em}#1\par\vspace{0.15em}}
\newcommand{\detailem}[1]{\noindent\hspace*{1.5em}\textit{#1}\par\vspace{0.15em}}

\begin{document}

\begin{center}
{\LARGE \textbf{Anderson Henrique da Silva}}\\[0.5em]
Center for the Study of the Metropolis (CEM)\\
Universidade de São Paulo --- USP Leste (EACH)\\
São Paulo, SP --- Brazil\\[0.3em]
\texttt{andersonheri@gmail.com} \\[0.15em]
\href{https://www.ahenriquecp.com}{www.ahenriquecp.com} \quad|\quad
\href{http://lattes.cnpq.br/7628365040600511}{Lattes} \quad|\quad
\href{https://orcid.org/0000-0002-1842-2725}{ORCID} \quad|\quad
\href{https://osf.io/ysj2q/}{OSF} \quad|\quad
\href{https://github.com/andersonheri}{GitHub}
\end{center}

\vspace{1em}

\section{Appointments}

\entry{Postdoctoral Research Fellow}{2025--present}
\detail{Center for the Study of the Metropolis (CEM), Universidade de São Paulo. FAPESP grant no.\ 2025/15250-0.}

\entry{Invited Researcher}{2025--present}
\detail{Institute for Applied Economic Research (IPEA), Brasília. Project: \textit{Intergovernmental Transfers in Brazilian Federalism} (PNPD).}

\entry{Research Assistant}{2024--2025}
\detail{Institute for Applied Economic Research (IPEA), Brasília.}

\entry{Collaborating Professor}{2023--2025}
\detail{Professional Master in Sociology in National Network (ProfSocio), Universidade Federal do Vale do São Francisco (UNIVASF).}

\entry{Substitute Professor of Political Science}{2022--2024}
\detail{Universidade Federal do Vale do São Francisco (UNIVASF), Juazeiro/BA.}

\section{Education}
\cvsubsection{Degree Programs}

\entry{Ph.D., Political Science}{2025}
\detail{Universidade Federal de Pernambuco (UFPE), Brazil.}
\detailem{Dissertation: ``Missão: Imprevisto --- Governança, Capacidades Estatais e Desempenho Municipal na Gestão de Risco e Desastres no Brasil.''}
\detail{Advisor: Dalson Britto Figueiredo Filho. Fields: State and Government; Local Power Studies; Intergovernmental Relations.}

\entry{M.A., Political Science}{2021}
\detail{Universidade Federal de Pernambuco (UFPE), Brazil.}
\detailem{Thesis: ``Desastres naturais, recursos federais de emergência e eleições municipais no Brasil (2012).''}
\detail{Advisor: Gabriela da Silva Tarouco. Co-advisor: Davi Cordeiro Moreira.}

\entry{B.A., Political Science}{2018}
\detail{Universidade Federal de Pernambuco (UFPE), Brazil. Advisor: Mariana Batista da Silva.}

\entry{Exchange, International Relations}{2016}
\detail{Universidade de Coimbra, Portugal. Advisor: Carmen Amado Mendes.}

\cvsubsection{Non-Degree Programs}

\entryplain{\textbf{R2}, Research Transparency and Reproducibility Training (32h) --- University of California, Berkeley.}{2024}

\entryplain{\textbf{Survey Data Analysis} (200h) --- University of Essex, United Kingdom.}{2019}

\entryplain{\textbf{Legislative Analysis Workshop} (6h) --- Universidade Federal de Minas Gerais (UFMG).}{2022}

\section{Grants and Fellowships}

\entryplain{FAPESP Postdoctoral Fellowship (proc.\ 2025/15250-0)}{2025--present}
\entryplain{FACEPE Doctoral Fellowship}{2021--2025}
\entryplain{FACEPE Master's Fellowship}{2019--2021}
\entryplain{CNPq Undergraduate Research Fellowship (PIBIC)}{2014--2017}
\entryplain{Santander Universities Ibero-American Scholarship}{2015}
\entryplain{Tutorial Education Program Scholarship (PET/MEC), UFPE}{2014--2015}

\section{Awards and Honors}

\entryplain{Ph.D.\ dissertation recommended for award by the examining committee, PPGCP/UFPE}{2025}
\entryplain{Research Transparency and Reproducibility Training (RT2) Scholarship, UC Berkeley}{2024}
\entryplain{TSE Summer School Scholarship, Toulouse School of Economics}{2023}
\entryplain{Honorable Mention, Institutional Politics Working Group (1st place), VII FBCP}{2022}
\entryplain{Third Best Paper Overall, VII Brazilian Forum of Graduate Studies in Political Science, UFMG}{2022}
\entryplain{Global South Academic Network Scholarship}{2019}
\entryplain{Vote of Applause, Recife City Council}{2017, 2019}

\section{Publications}
\cvsubsection{Peer-reviewed articles (accepted)}
\begin{enumerate}[leftmargin=0.3in,itemsep=0.4em,parsep=0pt]
\item Batista, M.; \textbf{Henrique, A.}; Porto, L. ``Investimento ou aprendizado? Incentivos da reeleição sobre a promoção de capacidades estatais nos municípios brasileiros.'' \textit{Revista Brasileira de Ciências Sociais} (forthcoming, 2026).
\item Palotti, P.\ L.\ M.; Linhares, P.\ T.\ F.; \textbf{Henrique, A.} ``Capacidades estatais municipais e o desempenho na gestão de transferências intergovernamentais da União.'' \textit{Texto para Discussão (IPEA)} (forthcoming, 2026).
\end{enumerate}

\cvsubsection{Peer-reviewed articles (published)}
\begin{enumerate}[leftmargin=0.3in,itemsep=0.4em,parsep=0pt]
\item \textbf{Henrique, A.}; Batista, M. ``Politização dos desastres naturais: alinhamento partidário, declarações de emergência e a alocação de recursos federais para os municípios no Brasil.'' \textit{Opinião Pública}, 26, 522--555, 2020.
\item Silva, D.\ P.\ A.; Figueiredo Filho, D.\ B.; \textbf{Henrique, A.} ``O poderoso NVivo: uma introdução a partir da análise de conteúdo.'' \textit{Política Hoje}, 24, 119--134, 2015.
\item Figueiredo Filho, D.\ B.; Santos Filho, R.\ P.; Rocha, E.\ C.; Silva Jr., J.\ A.; Silva, D.\ P.\ A.; \textbf{Henrique, A.}; Silva, L.\ E.\ O. ```Foi de morte matada': homicídios no Brasil em perspectiva comparada.'' \textit{Sistema Penal \& Violência}, 7, 6--17, 2015.
\item Figueiredo Filho, D.\ B.; Rocha, E.\ C.; Paranhos, R.; Silva Jr., J.\ A.; \textbf{Henrique, A.}; Silva, D.\ P.\ A.; Silva, L.\ E.\ O. ``Análise fatorial garantida ou o seu dinheiro de volta: uma introdução à redução de dados.'' \textit{Revista Eletrônica de Ciência Política --- RECP}, 5, 185--211, 2014.
\end{enumerate}

\cvsubsection{Book chapters}
\begin{enumerate}[leftmargin=0.3in,itemsep=0.4em,parsep=0pt]
\item \textbf{Henrique, A.}; Batista, M.; Porto, L.; Palotti, P.\ L.\ M. ``O Estado em números: construindo bons índices e interpretações de capacidade estatal.'' In: Gomide, A.; Marenco, A.; Oliveira, V.\ E.\ (Eds.), \textit{Reformar o Estado: o que fazer?}. Porto Alegre: Jacarta/ENAP, 2026, 224--238.
\item \textbf{Henrique, A.}; Soares, L.\ A.\ C. ``Governança do imprevisto: indicadores, diagnósticos e janelas de oportunidade na gestão de risco e desastres.'' In: Gomide, A.; Marenco, A.; Oliveira, V.\ E.\ (Eds.), \textit{Reformar o Estado: o que fazer?}. Porto Alegre: Jacarta/ENAP, 2026, 312--329.
\item \textbf{Henrique, A.}; Braga, D.\ D.\ A. ``Programa Brasil Alfabetizado: uma avaliação da teoria do programa.'' In: Nascimento, P.; Barros, A.\ T.\ D.\ L.\ (Eds.), \textit{Ciência Política: uma proposta educativa --- Políticas públicas, v.\ 2}. Campina Grande: Eduepb, 2023, 139--159.
\item Silva, A.\ H.; Batista, M. ``A politização dos desastres naturais: alinhamento partidário e a alocação de recursos federais de emergência para os municípios no Brasil.'' In: \textit{IV Simpósio de Direito Eleitoral do Nordeste}. Campina Grande: Abradep, 2022, 417--446.
\end{enumerate}

\cvsubsection{Op-eds}
\begin{enumerate}[leftmargin=0.3in,itemsep=0.4em,parsep=0pt]
\item Barros, A.\ T.\ D.\ L.; \textbf{Henrique, A.} ``Preocupémonos por la democracia en Brasil.'' \textit{El País}, 8 June 2022.
\item \textbf{Henrique, A.} ``O que há de político nos desastres naturais no Brasil?'' \textit{Estadão}, 19 February 2022.
\end{enumerate}

\section{Invited Talks and Presentations}

\entryplain{\textit{A política de gestão de risco e desastres no Brasil: quem, quando e como?} --- VI ENEPCP.}{2025}
\entryplain{\textit{Transparência e replicabilidade científica na escrita acadêmica.}}{2024}
\entryplain{\textit{Política e política pública da epidemia de Covid-19 no Brasil} (with Perich, R.; Figueiredo Filho, D.).}{2023}
\entryplain{\textit{Desastres naturais, recursos federais de emergência e eleições municipais} --- VII FBCP.}{2022}
\entryplain{\textit{Eleições no Brasil frente à crise democrática.}}{2022}
\entryplain{\textit{Teste de hipóteses e técnicas de comparação de médias aplicados à saúde.}}{2019}
\entryplain{\textit{O efeito do alinhamento político nas transferências federais de emergência no Brasil.}}{2019}
\entryplain{\textit{The map of public-private partnerships in sanitation in Brazil} --- 2nd LAPolMeth Meeting.}{2018}
\entryplain{\textit{More resources or a better local implementation?} --- VIII ALACIP.}{2015}

\section{Teaching Experience}

\cvsubsection{Graduate --- Professional Master in Sociology (ProfSocio, UNIVASF)}
\entryplain{Research Methodology}{2024, 2025}
\entryplain{MEDIALIT: Digital Media Literacy and Sociology Teaching in Schools}{2025}

\vspace{0.9em}
\cvsubsection{Undergraduate --- Social Sciences and Engineering (UNIVASF)}
\entryplain{Quantitative Methods and Statistical Interpretation}{2023}
\entryplain{Qualitative and Quantitative Methods}{2023}
\entryplain{Public Policy}{2022, 2023}
\entryplain{Political Theory I and III}{2022, 2023}
\entryplain{Federalism, Decentralization and Local Power}{2024}
\entryplain{Introduction to Scientific Work}{2024}
\entryplain{Sociology (for Civil, Production and Mechanical Engineering)}{2023, 2024}
\entryplain{Sociological Foundations of Psychology}{2024}
\entryplain{Philosophical Foundations of Psychology; Introduction to Philosophy}{2024}
\entryplain{Fundamentals of Anthropology and Sociology}{2022}

\vspace{0.9em}
\cvsubsection{Teaching Assistantships --- UFPE}
\entryplain{Public Policy I (with M.\ Batista)}{2022}
\entryplain{Research Seminar (with A.\ Steiner)}{2020}
\entryplain{Political Institutions I (with M.\ Batista)}{2019}
\entryplain{Anthropology (with P.\ Schörder)}{2018}
\entryplain{Quantitative Methods I; Political Theory I (with D.\ B.\ Figueiredo Filho)}{2015--2016}

\section{Research Groups and Networks}

\entryplain{Center for the Study of the Metropolis (CEM), USP --- Postdoctoral Fellow}{2025--present}
\entryplain{National Institute of Science and Technology on Quality of Government (INCT-QualiGov)}{2020--present}
\entryplain{Institutions, Government and Public Policy Research Group, DCP/UFPE}{2013--present}
\entryplain{Polikt --- Institutions, Participation and Political Culture, UNIVASF}{2022--present}
\entryplain{Local Crisis and Politics Consortium (CPL), UFABC/UFPR/UFPE}{2023--present}
\entryplain{DERN/GSAN Fellows' Network, University of Essex (Membership Team)}{2021--2022}

\section{Professional Service}
\cvsubsection{Editorial Board}
\entryplain{\textit{Policy Making --- Latin American Public Policy Review}}{2025--present}

\vspace{0.9em}
\cvsubsection{Ad hoc reviewer}
\textit{Opinião Pública}; \textit{Política \& Trabalho} (UFPB); \textit{Agenda Política}; \textit{Journal of Public Policy and Government}; \textit{Política Hoje} (UFPE); \textit{E-Legis --- Revista Eletrônica do PPG da Câmara dos Deputados}; ANPOCS Best Thesis and Dissertation Award (2026); SEBRAE Empreendedor Prefeitura Award (2024).

\section{Public Competitive Examinations (Approved)}

\entryplain{UNIRIO --- Political Science / Research Design. 3rd place overall; 1st in PPPIQ vacancy.}{2026}
\entryplain{UFAL --- Political Science. 3rd place.}{2025}
\entryplain{UERJ --- International Politics / IR Methods. 2nd place.}{2024}
\entryplain{UFCG --- Political Science. 5th place.}{2024}
\entryplain{UNIVASF --- Political Science (Substitute Professor). 1st place.}{2022}

\section{Skills and Languages}
\cvsubsection{Software}
R, Stata, SPSS, NVivo, \LaTeX, QGIS.

\cvsubsection{Languages}
Portuguese (native); English (fluent); Spanish (intermediate).

\vspace{1.2em}
\begin{center}
\footnotesize Last updated: \today
\end{center}

\end{document}
